Researchers \must justify all benchmark and metric choices in the \paper and \should summarize benchmark structure, task types, and limitations. They \should operationally define the phenomenon being measured, justify the sampling strategy used to select problems for inclusion in the benchmark, isolate the target capability from confounders where possible, and report an error analysis of the observed failures. Where possible, traditional (non-LLM) baselines \should be used for comparison. Researchers \must explain in the \paper why the selected metrics are suitable for the specific study; prior adoption in related work alone does not constitute sufficient justification. They \should report established metrics to make study results comparable, but can report additional metrics that they consider appropriate. Due to the inherent non-determinism of LLMs, experiments \should be repeated; the result distribution \should then be reported using descriptive statistics. If comparing models or tools, researchers \must use appropriate inferential statistics rather than relying solely on summary statistics. Researchers \should justify the number of experiment repetitions, for example through a power analysis or by monitoring convergence of descriptive statistics. Latency \must be reported when it can affect study outcomes (e.g., interactive user studies, latency comparisons).